
Parameter-efficient fine-tuning methods such as LoRA have been applied to State Space Models (SSMs) with variable success across different tasks. This work investigates the spectral dynamics of SSM adaptation by introducing the spectral memory horizon H\_spec, defined as -1/log|lambda\_max|, which characterizes the theoretical persistence length of information in the slowest-decaying eigenmode. Through experiments on Mamba-1.4B, the following findings are reported: (1) H\_spec is an input-independent property of pretrained models with coefficient of variation effectively zero (CV = 2.22 x $10^-16$ across 1000 sequences); (2) projection-only LoRA targeting in\_proj and x\_proj matrices preserves eigenvalues exactly (delta H\_spec = 0.0\%, correlation = 1.0); and (3) energy redistribution toward slow eigenmodes does not occur under projection-only LoRA (delta E = 5.93 x $10^-7$ nats, six orders of magnitude below the 0.1 nats threshold). This third finding, a negative result, eliminates the Eigenmode Utilization Hypothesis and provides evidence supporting the Memory Horizon Separation Hypothesis: that task adaptation via projection-only LoRA is bounded by the intrinsic spectral horizon. The discriminative task-failure test (H-M4) was not executed; thus, direct verification of task failure beyond H\_spec remains for future work.
